% Reviewer 1; OpenReview ID: gURK
% Original review text, followed by intentionally blank reply areas.
\section{Reviewer 1 (\texttt{gURK})}
\label{reviewer:1}
\noindent\textbf{Rating:} 5: Marginally below acceptance threshold\par
\noindent\textbf{Confidence:} 4: The reviewer is confident but not absolutely certain that the evaluation is correct

\subsection*{Paper Summary}
This paper starts from the argument that in OPD the student model's per token level agreement with teacher model is not perfect way for improving its reasoning process, as a student's generation can be moving towards the correct solution while having a high KL due to a different distribution from the teacher's. This paper then presents R\textsuperscript{2}-OPD, which uses a teacher-independent progress signal only as a reliability filter on the distillation loss. By splitting the student model's thinking trajectories into multiple segments by a preset markers, and process the reward with segments that contributes to the final ground truth answer, this method improves reasoning-oriented OPD, especially on harder benchmarks.

\subsection*{Strengths}
\begin{itemize}
\item This paper is well motivated and generally well written with a clear separation from distillation compatibility from reasoning progress, which targets the fact that OPD treats the former as a proxy for the latter and thus impacts the training efficiency.
\item The method is straightforward on the design including segmentation methods and easy to adopt for other methods. Specifically on the segmentation, the progress signal acts as a gate when applies the rules, and falls back to OPD when not available, the conflicts inside response are also addressed.
\item The experiments are well designed and clearly demonstrated, with Fig 1 showing student rollouts on DAPO that contain the answer but go unscored, which is a concrete argument that outcome-only filtering discards usable trajectories.On Fig 3-5 the paper also shows the configurations and ablations over segments merging strategy etc.
\end{itemize}

\subsection*{Weaknesses}
\begin{itemize}
\item This paper seems a combination of several existing works and some new components are not shown to be necessary through ablations like process reward or the process of filtering OPD supervision with the teacher-independent signal. It seems the genuine new components are mostly segmentation methods, conflict resolving rule and top-q masking, which are mostly heuristic and not compared with similar or simpler methods.
\par\Needspace{6\baselineskip}
\begin{responsebox}[frametitle={Response to W1}]
% Write the response here. All response content is intentionally blank.
\vspace*{2\baselineskip}
\end{responsebox}

\item The overall performance imprvements seem limited and specific to one benchmark family. In table 2 the gain is +0.43 avg@4, and over AIME25 avg@4 decreases.
\par\Needspace{6\baselineskip}
\begin{responsebox}[frametitle={Response to W2}]
% Write the response here. All response content is intentionally blank.
\vspace*{2\baselineskip}
\end{responsebox}

\item The evaluation design and results can be better explained. First of as the paper also suggests, only a small fraction of DAPO rollouts are scored, training at 7k tokens and evaluating at 8k might be truncating R1-Distill-1.5B responses as well, so shortening the responses will naturally improve the performance. Also it seems most baselines E-OPD, TIP-OPD, IW-OPD and Uni-OPD all score below the untrained student on AIME24, it might indicates the specific configuration favors the method in the paper. Moreover, the teacher model JustRL-DeepSeek-1.5B, is based on DeepSeek-R1-Distill-Qwen-1.5B, which comes with a low-divergence behavior and does not really reflect the motivation of this paper.
\par\Needspace{6\baselineskip}
\begin{responsebox}[frametitle={Response to W3}]
% Write the response here. All response content is intentionally blank.
\vspace*{2\baselineskip}
\end{responsebox}

\item The theory part is thin, Lemma 1 holds for fixed runs only which is natually too idealized. With prop.2, neither the 2 really analyzes the main contribution of the segmentation or masking strategy.
\par\Needspace{6\baselineskip}
\begin{responsebox}[frametitle={Response to W4}]
% Write the response here. All response content is intentionally blank.
\vspace*{2\baselineskip}
\end{responsebox}

\end{itemize}

\subsection*{Justification of Recommendation}
Based on the 3 weakness over limited performance improvements and evaluation design, I'm leaning towards a boarderline rejection.

\subsection*{Specific Points of Feedback for Rebuttal}
In addition to the weakness,

\begin{itemize}
\item The evaluation results could benefit from avg@32 on AIME with different seeds, more benchmarks.
\par\Needspace{6\baselineskip}
\begin{responsebox}[frametitle={Response to Q1}]
% Write the response here. All response content is intentionally blank.
\vspace*{2\baselineskip}
\end{responsebox}

\item Clear demonstration of segmentation and masking, e.g. the fraction of masked tokens that occur after the ground-truth answer first appears
\par\Needspace{6\baselineskip}
\begin{responsebox}[frametitle={Response to Q2}]
% Write the response here. All response content is intentionally blank.
\vspace*{2\baselineskip}
\end{responsebox}

\item Response length and truncation rate at 8k for student models over baselines
\par\Needspace{6\baselineskip}
\begin{responsebox}[frametitle={Response to Q3}]
% Write the response here. All response content is intentionally blank.
\vspace*{2\baselineskip}
\end{responsebox}

\item Justification of hyperparameters and implementation method of the baselines
\par\Needspace{6\baselineskip}
\begin{responsebox}[frametitle={Response to Q4}]
% Write the response here. All response content is intentionally blank.
\vspace*{2\baselineskip}
\end{responsebox}

\end{itemize}

\par\Needspace{12\baselineskip}
\subsection*{Final Response for Submission}
% Consolidate the final reply to this reviewer here before submission.
\par\Needspace{9\baselineskip}
\begin{responsebox}[frametitle={Final Response to Reviewer 1 (gURK)},linecolor=finalgreen,backgroundcolor=finalfill]
% Write the response here. All response content is intentionally blank.
\vspace*{4\baselineskip}
\end{responsebox}

